\documentclass[10pt,a4paper]{amsart}
\usepackage{amsmath,amssymb,amsthm,graphicx}
% Retained minimal preamble for the Stacks Project chapter flat.tex (More on Flatness).
% Verbatim excerpt of the chapter preamble: Xy-pic setup (native preamble lines 14-18) and the native
% theorem declarations with the mathematical macros (native preamble lines 154-217). Copied mechanically;
% nothing added, removed or redefined. The wrapper class and the amsmath/amssymb/amsthm/graphicx base
% packages are supplied by the assembler, so the class/externaldocument part of the native preamble is omitted.
% For commutative diagrams we use Xy-pic
\usepackage[all]{xy}

% We use 2cell for 2-commutative diagrams.
\xyoption{2cell}
% Theorem environments.
%
\theoremstyle{plain}
\newtheorem{theorem}[subsection]{Theorem}
\newtheorem{proposition}[subsection]{Proposition}
\newtheorem{lemma}[subsection]{Lemma}

\theoremstyle{definition}
\newtheorem{definition}[subsection]{Definition}
\newtheorem{example}[subsection]{Example}
\newtheorem{exercise}[subsection]{Exercise}
\newtheorem{situation}[subsection]{Situation}

\theoremstyle{remark}
\newtheorem{remark}[subsection]{Remark}
\newtheorem{remarks}[subsection]{Remarks}

\numberwithin{equation}{subsection}

% Macros
%
\def\lim{\mathop{\mathrm{lim}}\nolimits}
\def\colim{\mathop{\mathrm{colim}}\nolimits}
\def\Spec{\mathop{\mathrm{Spec}}}
\def\Hom{\mathop{\mathrm{Hom}}\nolimits}
\def\Ext{\mathop{\mathrm{Ext}}\nolimits}
\def\SheafHom{\mathop{\mathcal{H}\!\mathit{om}}\nolimits}
\def\SheafExt{\mathop{\mathcal{E}\!\mathit{xt}}\nolimits}
\def\Sch{\mathit{Sch}}
\def\Mor{\mathop{\mathrm{Mor}}\nolimits}
\def\Ob{\mathop{\mathrm{Ob}}\nolimits}
\def\Sh{\mathop{\mathit{Sh}}\nolimits}
\def\NL{\mathop{N\!L}\nolimits}
\def\CH{\mathop{\mathrm{CH}}\nolimits}
\def\proetale{{pro\text{-}\acute{e}tale}}
\def\etale{{\acute{e}tale}}
\def\QCoh{\mathit{QCoh}}
\def\Ker{\mathop{\mathrm{Ker}}}
\def\Im{\mathop{\mathrm{Im}}}
\def\Coker{\mathop{\mathrm{Coker}}}
\def\Coim{\mathop{\mathrm{Coim}}}

% Boxtimes
%
\DeclareMathSymbol{\boxtimes}{\mathbin}{AMSa}{"02}

%
% Macros for moduli stacks/spaces
%
\def\QCohstack{\mathcal{QC}\!\mathit{oh}}
\def\Cohstack{\mathcal{C}\!\mathit{oh}}
\def\Spacesstack{\mathcal{S}\!\mathit{paces}}
\def\Quotfunctor{\mathrm{Quot}}
\def\Hilbfunctor{\mathrm{Hilb}}
\def\Curvesstack{\mathcal{C}\!\mathit{urves}}
\def\Polarizedstack{\mathcal{P}\!\mathit{olarized}}
\def\Complexesstack{\mathcal{C}\!\mathit{omplexes}}
% \Pic is the operator that assigns to X its picard group, usage \Pic(X)
% \Picardstack_{X/B} denotes the Picard stack of X over B
% \Picardfunctor_{X/B} denotes the Picard functor of X over B
\def\Pic{\mathop{\mathrm{Pic}}\nolimits}
\def\Picardstack{\mathcal{P}\!\mathit{ic}}
\def\Picardfunctor{\mathrm{Pic}}
\def\Deformationcategory{\mathcal{D}\!\mathit{ef}}

% --- appended 2026-10-12 (Chief): the project own definitions of the reference/slogan/history
% display environments, copied mechanically from the shared Stacks preamble lines 8-12
% (\usepackage{verbatim} plus the three \newenvironment lines, which the chapter preamble keeps
% verbatim from the project preamble). Nothing added, removed or redefined: these environments
% carry no mathematical content and the project itself typesets them as comments. Without these
% three lines a counted statement whose own text uses them fails with Environment slogan undefined.
\usepackage{verbatim}
\newenvironment{reference}{\comment}{\endcomment}
\newenvironment{slogan}{\comment}{\endcomment}
\newenvironment{history}{\comment}{\endcomment}
\usepackage{cleveref}
\crefname{originallemma}{original lemma}{original lemmas}
\Crefname{originallemma}{Original lemma}{Original lemmas}
\makeatletter
\newlabel{algebra-lemma-composition-flat}{{00HC}{source}{}{sourcepointer-1}{}}
\newlabel{algebra-lemma-composition-flat@cref}{{[originallemma][0][]00HC}{[1][0][]source}{}{}{}}
\makeatother
\crefname{originallemma}{original lemma}{original lemmas}
\Crefname{originallemma}{Original lemma}{Original lemmas}
\makeatletter
\newlabel{morphisms-lemma-base-change-module-flat}{{01U8}{source}{}{sourcepointer-2}{}}
\newlabel{morphisms-lemma-base-change-module-flat@cref}{{[originallemma][0][]01U8}{[1][0][]source}{}{}{}}
\makeatother
\crefname{originallemma}{original lemma}{original lemmas}
\Crefname{originallemma}{Original lemma}{Original lemmas}
\makeatletter
\newlabel{morphisms-lemma-diagonal-unramified-morphism}{{02GE}{source}{}{sourcepointer-3}{}}
\newlabel{morphisms-lemma-diagonal-unramified-morphism@cref}{{[originallemma][0][]02GE}{[1][0][]source}{}{}{}}
\makeatother
\crefname{originallemma}{original lemma}{original lemmas}
\Crefname{originallemma}{Original lemma}{Original lemmas}
\makeatletter
\newlabel{morphisms-lemma-etale-flat}{{02GS}{source}{}{sourcepointer-4}{}}
\newlabel{morphisms-lemma-etale-flat@cref}{{[originallemma][0][]02GS}{[1][0][]source}{}{}{}}
\makeatother
\crefname{originallemma}{original lemma}{original lemmas}
\Crefname{originallemma}{Original lemma}{Original lemmas}
\makeatletter
\newlabel{morphisms-lemma-etale-smooth-unramified}{{02GK}{source}{}{sourcepointer-5}{}}
\newlabel{morphisms-lemma-etale-smooth-unramified@cref}{{[originallemma][0][]02GK}{[1][0][]source}{}{}{}}
\makeatother
\title{Flatness is the same over the base and over an etale cover}
\author{The Stacks Project Authors}
\date{stacks-project:flat.tex:local-snapshot-20261011; GNU FDL 1.2 or later}
\begin{document}
\maketitle
\noindent\textit{Editorial note.} All delivered mathematical text is the author's own, copied from the pinned Stacks Project chapter file flat.tex; no statement or proof step was rewritten, reordered or removed.\par\medskip
\noindent\textit{Editorial note.} Preamble provenance: the chapter's own retained minimal preamble (sources/stacks-flat/retained-preamble.tex, lines 1-85) is a verbatim excerpt of the native Xy-pic setup and of the native theorem declarations and mathematical macros, followed by the project's own reference/slogan/history display-environment definitions copied mechanically from the shared preamble; the wrapper class and the amsmath/amssymb/amsthm/graphicx base packages are supplied by the assembler, and extra\_packages is left empty because everything needed is already in the retained preamble. No macro was invented or redefined.\par\medskip
\noindent\textit{Editorial note.} Same-chapter pointers are located at their \textbackslash{}label lines inside the pinned chapter file; cross-chapter pointers are located in the project's own tag list staged in the snapshot (sources/stacks-chow/tags.txt), at the verbatim lines carrying their canonical tags.\par\medskip
\noindent\textit{Allowed standard prerequisites (editorial).}\begin{itemize}
\item An etale morphism is flat, and flatness is preserved by composition. Morphisms, label morphisms-lemma-etale-flat (proof line 139); Algebra, label algebra-lemma-composition-flat (proof line 141). Gives the easy implication from flatness over the etale scheme to flatness over the base.
\item Base change of a flat module, and the diagonal of an etale morphism is an open immersion. Morphisms, labels morphisms-lemma-base-change-module-flat (proof line 147), morphisms-lemma-diagonal-unramified-morphism and morphisms-lemma-etale-smooth-unramified (proof lines 150-151). The first makes the pullback of the module to the product flat over the etale scheme, and the others identify the source with an open subscheme of the product on which the restriction of the second projection is the given morphism.
\end{itemize}\medskip
\ifdefined\hypertarget\hypertarget{sourcepointer-1}{}\fi
\noindent\textit{Source reference.} Original lemma 00HC: sources/stacks-chow/tags.txt, source lines 638--638. Navigation pointer for \textbackslash{}ref\{algebra-lemma-composition-flat\} used inside the counted span (line 141); the argument that consumes it is retained in the span.\par\medskip
\ifdefined\hypertarget\hypertarget{sourcepointer-2}{}\fi
\noindent\textit{Source reference.} Original lemma 01U8: sources/stacks-chow/tags.txt, source lines 2401--2401. Navigation pointer for \textbackslash{}ref\{morphisms-lemma-base-change-module-flat\} used inside the counted span (line 147); the argument that consumes it is retained in the span.\par\medskip
\ifdefined\hypertarget\hypertarget{sourcepointer-3}{}\fi
\noindent\textit{Source reference.} Original lemma 02GE: sources/stacks-chow/tags.txt, source lines 3225--3225. Navigation pointer for \textbackslash{}ref\{morphisms-lemma-diagonal-unramified-morphism\} used inside the counted span (line 150); the argument that consumes it is retained in the span.\par\medskip
\ifdefined\hypertarget\hypertarget{sourcepointer-4}{}\fi
\noindent\textit{Source reference.} Original lemma 02GS: sources/stacks-chow/tags.txt, source lines 3239--3239. Navigation pointer for \textbackslash{}ref\{morphisms-lemma-etale-flat\} used inside the counted span (line 139); the argument that consumes it is retained in the span.\par\medskip
\ifdefined\hypertarget\hypertarget{sourcepointer-5}{}\fi
\noindent\textit{Source reference.} Original lemma 02GK: sources/stacks-chow/tags.txt, source lines 3231--3231. Navigation pointer for \textbackslash{}ref\{morphisms-lemma-etale-smooth-unramified\} used inside the counted span (line 151); the argument that consumes it is retained in the span.\par\medskip
\begin{lemma}
\label{lemma-etale-flat-up-down}
Let $X \to T \to S$ be morphisms of schemes with $T \to S$ \'etale.
Let $\mathcal{F}$ be a quasi-coherent $\mathcal{O}_X$-module.
Let $x \in X$ be a point. Then
$$
\mathcal{F}\text{ flat over }S\text{ at }x
\Leftrightarrow
\mathcal{F}\text{ flat over }T\text{ at }x
$$
In particular $\mathcal{F}$ is flat over $S$ if and only if $\mathcal{F}$
is flat over $T$.
\end{lemma}

\begin{proof}
As an \'etale morphism is a flat morphism (see
Morphisms, Lemma \ref{morphisms-lemma-etale-flat})
the implication ``$\Leftarrow$'' follows from
Algebra, Lemma \ref{algebra-lemma-composition-flat}.
For the converse assume that $\mathcal{F}$ is flat at $x$ over $S$.
Denote $\tilde x \in X \times_S T$ the point lying over $x$ in $X$
and over the image of $x$ in $T$ in $T$.
Then $(X \times_S T \to X)^*\mathcal{F}$ is flat at $\tilde x$ over $T$
via $\text{pr}_2 : X \times_S T \to T$, see
Morphisms, Lemma \ref{morphisms-lemma-base-change-module-flat}.
The diagonal $\Delta_{T/S} : T \to T \times_S T$ is an open immersion;
combine
Morphisms, Lemmas \ref{morphisms-lemma-diagonal-unramified-morphism} and
\ref{morphisms-lemma-etale-smooth-unramified}.
So $X$ is identified with open subscheme of $X \times_S T$,
the restriction of $\text{pr}_2$ to this open is the given morphism $X \to T$,
the point $\tilde x$ corresponds to the point $x$ in this open, and
$(X \times_S T \to X)^*\mathcal{F}$ restricted to this open is $\mathcal{F}$.
Whence we see that $\mathcal{F}$ is flat at $x$ over $T$.
\end{proof}

\end{document}
